\documentclass[a4paper]{article}
%To  colorize table
\usepackage[table]{xcolor}
\usepackage[a4paper, total={7in, 11.0in},includefoot, includehead, headsep=24pt, headheight=4cm]{geometry}
\usepackage{layout}%No need for this package at production.
\usepackage{setspace}%For switch between line spacing.
\usepackage{graphicx}
\usepackage{multirow}
%For drawing lines and arbitrary shapes:
\usepackage{tikz}
%To set fixed length for table: 
\usepackage{array}
\usepackage{ifthen}
%To draw squares:
\usepackage{amssymb}
%To trim strings:
\usepackage{trimspaces}
\usepackage{ifthen}
%To draw dynamic multi-column table:
\usepackage{etoolbox}
\usepackage{float}

%\usepackage{biblatex}
%\addbibresource{refs.bib} % with extension
%\setcitestyle{authoryear,open={(},close={)}}
%\bibliographystyle{ieeetr-fa.bst}
%For references:
% \usepackage[style=plain]{biblatex}
%\bibliographystyle{ieeetr}



% \usepackage{biblatex}
% \addbibresource{references.bib}
%\usepackage[style=numeric]{biblatex}




%\renewcommand{\bibsection}{}{\rl{{مراجع}\hfill}} 


%\addbibresource{references.bib}


%List of needed counters:
\newcounter{length}
\newcounter{itemCount}


\title{Proposal}
\author{Name}
\date{\today}

\newcommand{\blankField}{ $\ldots\ldots\ldots\ldots\ldots$}
\makeatletter
\newcommand{\trim}[1]{\trim@spaces@noexp{#1}}
\makeatother
\newcolumntype{P}[1]{>{\centering\arraybackslash}p{#1}}


\usepackage{fancyhdr}
\pagestyle{fancy}
\fancyhf{}
\usepackage{xepersian}
\settextfont[Scale=1.2]{Yas}
\setlatintextfont[Scale=1]{Times New Roman}
%\bibliographystyle{ieeetr-fa.bst}
%\bibliographystyle{your-bst-file} % Replace 'your-bst-file' with the actual BST file name (without the .bst extension)

\bibliography{references} 
\rhead{
    \begin{tabular}{c}
         \includegraphics[width=3cm, height=3cm]{./logo-fa-IR.png}\\
    \end{tabular}
}
\chead{
    \setstretch{2}
    بسمه تعالی \\
    دانشکده مهندسی هوافضا \\
    \textbf{\large فرم تعریف پروژه کارشناسی ارشد}
}

\lhead{
    \ifthenelse{\isodd{\value{page}}}
        {
            \setstretch{1.5}
            \begin{tabular}{c c}
                تاریخ: &
                 1402/08/10 \\
                شماره: & \blankField\\
                پیوست: & \blankField                
                % \blankField & تاریخ: \\
                % \blankField & شماره:\\
                % \blankField & پیوست:
            \end{tabular}
        }
        {
            \setstretch{1.5}
            \begin{tabular}{c c}
                تاریخ:
                 & 1402/08/10 \\
                شماره: & \blankField\\
                پیوست: & \blankField
                
                % \blankField & تاریخ: \\
                % \blankField & شماره:\\
                % \blankField & پیوست:
            \end{tabular}
        }
}
% cite

%Improvement: take advantages of the optional parameters! Currently I don't have any time to learn and implement it.
\newcommand{\Information}[9]{
{
    \setstretch{1.5}
    \small
    \bfseries
    \begin{tabular}{r r r}
         نام و نام خانوادگی : #1
         &
         شماره دانشجویی: #2
         &
         معدل: #3
         \\
         گرایش: #4
         &
         تعداد واحدهای گذرانده: #5
         &
         استاد راهنما: #6
         \\
         استاد راهنمای همکار: #7
         &
         تعداد واحد پروژه: #8
         &
         استاد ممتحن: #9
    \end{tabular}
}
}
\usepackage{array} % Add this package to use 'm' column specifier

\newcommand{\Title}[3]{
	{
		\setstretch{1.5}
		\bfseries
		\noindent
		عنوان کامل پروژه:
		
		فارسی:
        
        #1
		
		انگلیسی:


		\lr{#2}\\
		
		\noindent
		نوع پروژه:\;
		\begin{tabular}{m{10em} m{10em} m{10em}}
			کاربردی:
			\ifthenelse{#3 = 0 \or #3 = 3 \or #3 = 4}{
				$\blacksquare$
			}{
				$\square$
			}
			&
			بنیادی:
			\ifthenelse{#3 = 1 \or #3 = 3}{
				$\blacksquare$
			}{
				$\square$
			}
			&
			توسعه‌ای:
			\ifthenelse{#3 = 2 \or #3 = 3 \or #3 = 4}{
				$\blacksquare$
			}{
				$\square$
			} \\
		\end{tabular}
	}
}

% \newcounter{length}
% \newcounter{itemCount}

\newcommand{\Description}[4]{
    \setstretch{1.5}
    \noindent
    \textbf{#1:}
    \noindent
    
    #2

    \vspace{20pt}
    
    \setcounter{length}{0}
    \setcounter{itemCount}{0}
    %First obtain the length of field
    \foreach\x in#4{%
        \addtocounter{length}{1}
    }
    
    \ifthenelse{#3 = 1}{
    \textbf{کلمات کلیدی:}\;
    \foreach\x in#4{%
        \addtocounter{itemCount}{1}
        \trim{\x}
        \ifthenelse{\thelength = \theitemCount}{}{-}
    }
    \vspace{1em}
    }
}


\makeatletter
\newcommand{\ProgressTable}[1]{
{
    \noindent
    \textbf{مراحل انجام پروژه و زمان‌بندی آن:}
    \\
    
    \def\tableData{}
    
    \foreach \x\y\z in #1{
        \protected@xappto\tabledata{ \x & \y & \z\\ \noexpand\hline}
    }
    
    %Add space to each row in table:
    \renewcommand{\arraystretch}{1.5}
    \begin{center}
    \begin{tabular}{| c | m{0.8\textwidth} | c |}
        \hline
        \tabledata
    \end{tabular}
    %Add space to each row in table:
        
    \end{center}
    \renewcommand{\arraystretch}{1}
}
    
}
\makeatother

%\newcommand{\References}{
%{
%    \noindent
%    \nocitep*}
%    \textbf{مراجع:}
%    \begin{latin}
%        \printbibliography[heading=none]
%    \end{latin}
%}
%}
\newcommand{\Courses}[1]{
{
    \noindent
    \textbf{دروس مورد نیاز:}\\
    
    %Not completed yet.
    \centering
    %Add space to each row in table:
    \renewcommand{\arraystretch}{1.5}
    \begin{center}
    \begin{tabular}{| >{\centering}m{6em} | >{\centering}m{3em} | >{\centering}m{6em} | >{\centering}m{6em} | >{\centering}m{3em} | >{\centering\arraybackslash}m{6em} |}
         \hline
         \rowcolor{lightgray}
         \multicolumn{3}{|c|}{جبرانی}
         &
         \multicolumn{3}{|c|}{تخصصی
            \tiny
            (ارتباط موضوع پروژه با دروسی که دانشجو گذرانده یا باید بگذراند)
         }
         \\
         \hline
         \rowcolor{lightgray}
         گذرانده
         &
         نمره
         &
         باید بگذراند
         &
         گذرانده
         &
         نمره
         &
         باید بگذراند
         \\
         \hline
         &
         &
         &
         &
         &
         \\
         \hline
         &
         &
         &
         &
         &
         \\
         \hline
         &
         &
         &
         &
         &
         \\
         \hline
    \end{tabular}
    \end{center}
    %Add space to each row in table:
    \renewcommand{\arraystretch}{1}
}
}
\newcommand{\Signatures}{
{
    %Add space to each row in table:
    \renewcommand{\arraystretch}{1.5}
    
    \begin{center}
        \scriptsize
        \begin{tabular}{|m{15em}|m{15em}|m{15em}|}
            \hline
            استاد راهنما:
            &
            نظر گروه:
            &
            نظر کمیته تحصیلات تکمیلی دانشکده:
            \\
            تاریخ تحویل فرم به مدیر گروه:
            &
            &
            \\
            امضای استاد راهنما:
            &
            تاریخ جلسه گروه:
            &
            تاریخ جلسه کمیته:
            \\ 
            &
            امضای مدیر گروه:
            &
            امضای معاون تحصیلات تکمیلی:
            \\
            &
            &
            \\
            &
            &
            \\
            & &\\
            & &\\
            \hline
        \end{tabular}
    \end{center}
    \footnotesize
    توجه: فرم تعریف پروژه بایستی یک روز قبل از جلسه گروه توسط استاد راهنما تحویل مدیر گروه شود.
    %Add space to each row in table:
    \renewcommand{\arraystretch}{1}
}
}

%\renewcommand{\headrulewidth}{0pt}
%\renewcommand{\footrulewidth}{0pt}


\begin{document}

%First parameter: Name of student(Your name)
%Second parameter: Student Number
%Third parameter: Student's GPA
%Fourth parameter: Student's Major
%Fifth parameter: Student's passed courses credits.
%Sixth parameter: Project Supervisor.
%Seventh parameter: Project Co-Supervisor
%Eighth parameter: Project Credits.
%Ninth parameter: Project Examiner
\Information{
    علی بنی‌اسد
}{
    401209244
}{
    16/34
}{
    فضا
}{
    15
}{
    دکتر نوبهاری
}{}{
    6
}{
%    نام استاد ممتحن
}

\noindent

%First parameter: Title of your project in Persian.
%Second parameter: Title of your project in English.
%Third parameter: Type of your project =>
%   0: Theoretical.
%   1: Practical.
%   2: Both practical & theoretical
\Title{
    هدایت  یادگیری تقویتی مقاوم مبتنی بر بازی دیفرانسیلی در محیط‌های پویای چندجسمی با پیشران کم}
    {
    Robust Reinforcement Learning Differential Game Guidance in Low-Thrust, Multi-Body Dynamical Environments
}{
    4
}

\noindent

\Description{معرفی موضوع}{
\noindent
 در سال‌های اخیر، پیشرفت‌های فناوری در زمینه‌های مختلف، از جمله کنترل پرواز، پردازش سیگنال و هوش مصنوعی، به افزایش کاربردهای ماهواره با پیشران کم در منظومه زمین‐ماه کمک کرده است. ماهواره با پیشران کم می‌تواند برای تعقیب ماهواره‌ها، انتقال مداری و استقرار ماهواره‌ها استفاده شود.
روش‌های هدایت بهینه قدیمی جهت کنترل ماهواره‌ها اغلب نیازمند فرضیات ساده کننده، منابع محاسباتی فراوان و شرایط اولیه مناسب هستند. الگوریتم‌های مبتنی بر یادگیری تقویتی این توانایی را دارند که بدون مشکلات اشاره‌شده هدایت ماهواره را انجام دهند. به همین دلیل، این الگوریتم‌ها می‌توانند امکان محاسبات درونی 
(\lr{On-board Computing})  
را فراهم می‌کنند.



%این فناوری می تواند به کاهش هزینه‌ها، بهبود عملکرد و کاهش خطرات مرتبط با پروازهای فضایی کمک کند.
%این پژوهش یک منطق هدایت مقاوم با استفاده از یادگیری تقویتی بازی‌های دیفرانسیلی 
%\LTRfootnote{Differential Game Reinforcement Learning (DGRL)}
% را ارائه می‌دهد. این منطق هدایت در بستر شبکه عصبی
% \LTRfootnote{Neural Network (NN)}
%  برای راهنمایی خودکار و حلقه بسته فضاپیماهایی با پیشران کم در محیط‌های چند جسمی متغیر مانند منظومه زمین-ماه، ارائه می‌شود.
  
}{0}{
    {
        کلمه کلیدی نمونه
        ,
        کلمه کلیدی نمونه
        ,
        کلمه کلیدی نمونه
    }
}

\Description{اهمیت موضوع}{
%مطالعه و توسعه این الگوریتم هدایت جدید با استفاده از الگوریتم‌های یادگیری تقویتی و بازی‌های دیفرانسیلی دارای اهمیت ویژه‌ای است. این الگوریتم هدایت، به عنوان یک مدل مقاوم، قابلیت مقابله با نوسانات و متغیرهای محیطی را داراست. علاوه بر این، ساختار جدید پردازنده‌های کامپیوتر پرواز، با امکانات به‌روز و محاسبات ماتریسی بهبود یافته، اجرای بهینه‌تر الگوریتم‌های یادگیری ماتریسی را آسان کرده و امکان اجرای موثر آنها را فراهم کرده است.
%
%با توجه به اینکه ایستگاه زمینی در بیشتر زمان ماموریت فاصله زمانی قابل توجهی از محل اجرای دستورات و مکان اعمال دستورات دارد، استفاده از این الگوریتم هدایت می‌تواند به عنوان یک تکنولوژی بسیار مهم در کاهش تاخیرهای ناشی از این فواصل زمانی و افزایش پایداری و دقت سیستم‌های کامپیوتر پرواز تلقی شود. توانایی این الگوریتم هدایت در انجام محاسبات پیچیده و تصمیم‌گیری‌های برنامه‌ریزی در زمان واقعی، اهمیت ویژه‌ای برای امنیت و پایداری پروازهای فضایی دارد و می‌تواند تاثیر مثبتی در اجرای موثر ماموریت‌های فضایی آینده داشته باشد.

%این الگوریتم هدایت از الگوریتم‌های یادگیری تقویتی و بازی‌های دیفرانسیلی استفاده می‌کند.
%این الگوریتم هدایت دارای مزایای مختلفی از جمله مقاومت در برابر نوسانات و متغیرهای محیطی، امکان اجرای بهینه‌تر الگوریتم‌های یادگیری ماتریسی، کاهش تاخیرهای ناشی از فواصل زمانی بین ایستگاه زمینی و فضاپیما، افزایش پایداری و دقت سیستم‌های کامپیوتر پرواز، و توانایی انجام محاسبات پیچیده و تصمیم‌گیری‌های برنامه‌ریزی در زمان واقعی است.
%این الگوریتم هدایت می‌تواند برای کاربردهای مختلفی در ماموریت‌های فضایی آینده مورد استفاده قرار گیرد.
%محاسبات درونی برای هدایت فضاپیما در محیط‌های پیچیده دینامیکی چند بدنه، از اهمیت بالایی برخوردار است. 
رویکردهای قدیمی برای هدایت فضاپیما در محیط‌های پیچیده چندجسمی اغلب به ایستگاه‌های کنترل زمینی متکی هستند که می‌توانند در معرض خرابی ارتباطات، تاخیرهای زمانی و محدودیت‌های انتقال داده باشند. لذا، ارائه روش‌های هدایت بهینه و مقاوم که امکان محاسبه فرامین هدایت را درون ماهواره فراهم کند، از اهمیت بالایی برخوردار است.
}{0}{{}}

\Description{کاربردها}{
% کاربردهای الگوریتم هدایت به‌صورت محاسبات درونی، در  ماموریت‌های مختلف فضایی در ادامه آورده شده است. این الگوریتم هدایت می‌تواند در ماموریت‌های فضایی مختلفی مانند ماموریت‌های ماهواره‌ای، ماموریت‌های سفر به ماه و ماموریت‌های نزدیک به زمین مورد استفاده قرار گیرد.
% الگوریتم هدایتی که مقاوم، بهینه و محاسبات آن به‌صورت درونی است توانایی انجام مانورها، تعقیب مسیر‌های مرجع بین مدارهای مختلف و بهبود اشتباهات را در محیط‌های پویا و پیچیده فضایی به‌صورت خودکار و بدون نیاز به ایستگاه زمینی را دارد. این نوع الگوریتم توانایی حل مسائل پیچیده هدایت فضایی که نیازمند تصمیم‌گیری در زمان واقعی و مطابق با متغیرهای محیطی می‌باشند، استفاده می‌شود. از جمله این مسائل می‌توان به تغییر مسیر در محیط‌های پیچیده، تنظیمات نیروی پیشرانه و اصلاح وضعیت نسبت به مدارهای مرجع اشاره کرد. از کاربردهای آن می‌توان به استفاده در ماموریت‌های مختلف فضایی را دارد. از جمله ماموریت‌های مختلف می‌توان به ماموریت‌های ماهواره‌ای، تعمیر و نگهداری ماموریت فضایی، تعقیب اهداف در مدارهای مختلف و انجام مانور‌های پیچیده اشاره کرد.
الگوریتم‌های هدایت مقاوم که محاسبات آن‌ها به‌صورت درون‌مداری انجام می‌شود، در ماموریت‌های فضایی مختلفی مانند ماهواره‌ای، سفر به ماه و نزدیک به زمین کاربرد دارند. این الگوریتم‌ها دارای مزایایی مانند استقلال از ایستگاه زمینی، سرعت تصمیم‌گیری و بهبود ایمنی هستند. از کاربردهای این الگوریتم‌ها می‌توان به کنترل مدار و وضعیت ماهواره‌ها، کنترل مدار و وضعیت فضاپیماها در طول سفر به ماه، کنترل مدار و وضعیت فضاپیماها در مدارهای نزدیک به زمین، تعمیر و نگهداری ماموریت‌های فضایی، تعقیب اهداف در مدارهای مختلف و انجام مانور‌های پیچیده اشاره کرد.

%\begin{enumerate}
%    \item \textbf{هدایت خودکار فضاپیماها:} 
%     الگوریتم هدایت اشاره شده توانایی انجام مانورها، تعقیب مسیر‌های مرجع بین مدارهای مختلف و بهبود اشتباهات را در محیط‌های پویا و پیچیده فضایی به‌صورت خودکار و بدون نیاز به ایستگاه زمینی را دارد.
%
%    \item \textbf{حل مسائل پیچیده هدایت:}  الگوریتم هدایت  اشاره شده برای حل مسائل پیچیده هدایت فضایی که نیازمند تصمیم‌گیری در زمان واقعی و مطابق با متغیرهای محیطی می‌باشند، استفاده می‌شود. از جمله این مسائل می‌توان به تغییر مسیر در محیط‌های پیچیده، تنظیمات نیروی پیشرانه و اصلاح وضعیت نسبت به مدارهای مرجع اشاره کرد.
%
%    \item \textbf{هدایت در محیط‌های نامطمئن:}
%     الگوریتم توانایی هدایت در محیط‌های نامطمئن و بدون نیاز به مدل دقیق محیط فضایی را دارد. این موضوع در مواقعی که مدل‌سازی دقیق مدل‌های دینامیکی مشکل باشد، قابل اجرا است.
%
%    \item \textbf{استفاده در ماموریت‌های متنوع:} این الگوریتم امکان استفاده در ماموریت‌های مختلف فضایی را دارد. از جمله ماموریت‌های مختلف می‌توان به ماموریت‌های ماهواره‌ای، تعمیر و نگهداری ماموریت فضایی، تعقیب اهداف در مدارهای مختلف و انجام مانور‌های پیچیده اشاره کرد.
%\end{enumerate}
}{0}{{}}


\Description{تعریف دقیق مسئله}{
	در پژوهش حاضر، مسأله هدایت فضاپیما برای سیستم زمین-ماه با استفاده از مدل دینامیکی سه‌جسم محدود دایره‌ای \lr{Circular Restricted Three-Body Problem (CR3BP)} مطرح می‌شود. در این مسأله، یک فضاپیما با سیستم پیشران کم به یک مسیر مرجع در سیستم \lr{CR3BP} منتقل می‌شود. شرایط اولیه فضاپیما از مسیر مرجع انحراف‌های تصادفی دارند. هدف از این پژوهش، توسعه یک الگوریتم هدایت حلقه‌بسته با استفاده از اصول یادگیری تقویتی بر مبنای بازی دیفرانسیلی است که به فضاپیما اجازه می‌دهد با فرض بدترین اغتشاش به مسیر مرجع بازگردد، مسیر مرجع را دنبال کرده و به مدار مقصد برسد. بازی دیفرانسیلی موجب می‌شود الگوریتم هدایت نه تنها بهینه، بلکه مقاوم نیز باشد.
}{0}{{}}


\Description{فرضیات مسئله}{
\begin{itemize}
	\item مدل دینامیکی مسأله، مسأله سه ‌جسم محدود دایره‌ای
	سیستم زمین-ماه است. این فرضیه یک سیستم دو جسمی ساده است که فضاپیما را به‌عنوان جسم سوم بسیار کوچک فرض می‌کند.
	\item حرکت فضاپیما تنها در صفحه اتفاق می‌افتد و تنها نیروهای گرانشی و اغتشاش در‌نظر گرفته می‌شوند.
	\item فضاپیما به یک سیستم پیشران کم مجهز است که شعاع تاثیر مشخصی دارد. میزان و جهت پیشران به‌صورت مستمر تغییر کرده و قابل تنظیم است.
	\item مسأله هدایت شامل انتقال بین مدارهای دوره‌ای (مدارهای لیاپانوف) در مدل \lr{CR3BP} زمین-ماه است.
%	\item  توانایی موتور دارای محدودیت است. موتور قادر به تولید پیشران در هر جهت و مقدار نیست.
%	\item وضعیت فضاپیما به طور کامل قابل مشاهده و در دسترس الگوریتم هدایت در هر مرحله زمانی است.
	\item محیط یادگیری تقویتی ویژگی‌های مارکوف را دارد؛ به این معنا که حالت فعلی تمام اطلاعات لازم برای پیش‌بینی حالت‌های آینده را فراهم می‌کند.
%	\item آموزش در یک محیط چرخه‌ای \LTRfootnote{episodic}
%	 انجام می‌شود.
	\item فضاپیما در مسیر نزدیک به مرجع با فرض بدترین اغتشاش می‌ماند تا به مدار مقصد برسد.
%	\item هدف اصلی تنها  اجرا پذیری مسأله   نیست و دیگر اهداف مانند مصرف کمتر سوخت یا زمان انتقال با استفاده از نظریه بازی در نظر گرفته می‌شوند.
%	\item این رویکرد نیاز به دسترسی به منابع محاسباتی با عملکرد بالا برای آموزش یادگیری تقویتی قبل از اجرای آن در رایانه پرواز فضاپیما را دارد.
\end{itemize}

	
}{0}{{}}
\newpage
\Description{روش انجام کار}{
	% مرحله 1: ایجاد محیط دینامیکی
% مراحل کلی
% مراحل کلی
\begin{itemize}
	\item \textbf{مرحله 1: مدل‌سازی و شبیه‌سازی محیط دینامیکی}
	\begin{itemize}
		\item پیاده‌سازی مدل دینامیکی محیط مورد نظر انجام می‌شود. این مدل باید شامل تمام جوانب مهم محیط از جمله قوانین حرکت، ابعاد فضایی و حالت اولیه باشد.
%		\item در صورت پیچیدگی مدل دینامیکی محیط، از کتابخانه‌ها و ابزارهای مهندسی نرم‌افزاری مانند \lr{NumPy} و \lr{SciPy} در پایتون برای پیاده‌سازی آن استفاده می‌شود.
		\item مدل دینامیکی محیط  به‌صورتی که اطلاعات مورد نیاز برای یادگیری تقویتی (مانند حالت فعلی) را تامین کند، ایجاد می‌شود.
	\end{itemize}
	
	\item \textbf{مرحله 2: ایجاد و آموزش بازیگر هدایت}
	\begin{itemize}
		\item در این مرحله، بازیگری برای انجام هدایت سفینه در محیط  ایجاد می‌شود. این بازیگر مسئول انتخاب اعمال (اعمال کم‌پیشران) بر اساس حالت فعلی محیط است.
		\item از الگوریتم‌های یادگیری تقویتی مختلف مانند \lr{Q-Learning}، \lr{Deep Q-Networks (DQN)}، یا \lr{Proximal Policy Optimization (PPO)} برای ایجاد بازیگر هدایت استفاده می‌شود. انتخاب الگوریتم بستگی به  محیط و اهداف آموزش دارد.
	\end{itemize}
	
%	\item \textbf{مرحله 3: استفاده از منابع \lr{GPU} در \lr{Google Colab}}
%	\begin{itemize}
%		\item \lr{Google Colab} به عنوان یک محیط آموزش آنلاین به کار می‌رود که اجازه می‌دهد کدهای پایتون در یک محیط مبتنی بر مرورگر اجرا شوند. از این مزیت این محیط می‌توان به دسترسی به منابع \lr{GPU} برای آموزش مدل‌های عمیق اشاره کرد.
%		\item برای استفاده از \lr{GPU} در \lr{Google Colab}، می‌توان از کتابخانه‌های معروفی مانند \lr{TensorFlow} یا \lr{PyTorch} که از \lr{GPU} پشتیبانی می‌کنند، استفاده می‌شود. به این ترتیب، آموزش مدل‌های یادگیری تقویتی بسیار سریع‌تر و رایگان انجام می‌شود.
%	\end{itemize}
\end{itemize}

	
}{0}{{}}
\vspace{-1cm}
\Description{پیشینه‌ی موضوع}{
	پژوهش‌های خارجی: 
	
	
%	رایج‌ترین تکنیک‌های هدایت معمولاً با استفاده از ایستگاه‌های کنترل مستقر در زمین انجام می‌شوند. با این حال، حساسیت به شکست ارتباطات، تاخیرهای زمانی، محدودیت‌های انتقال داده، پیچیدگی سنسورها و هزینه‌های عملیاتی همگی از دلایلی هستند که هدایت، مسیریابی و کنترل\LTRfootnote{Guidance, Navigation and Control (GNC)}
%	 را از کامپیوتر ایستگاه زمینی به کامپیوتر پرواز منتقل می‌کنند. 
%در حالی که در طراحی مسیر از افزایش قدرت سخت‌افزارهای کامپیوتر بهره‌برداری می‌کنند، تعداد کمی از آن‌ها برای اجرای خودکار درون‌سفینه عملی هستند. این به دلیل محدودیت‌های منابع محاسباتی است که در فضاپیماها وجود دارد.


هدایت فضاپیماها معمولاً با استفاده از ایستگاه‌های زمینی
% هدایت، مسیریابی و کنترل\LTRfootnote{Guidance, Navigation and Control (GNC)}  
 انجام می‌شود. با این حال، این تکنیک‌ها دارای محدودیت‌هایی از جمله حساسیت به قطع ارتباطات، تاخیرهای زمانی، و محدودیت‌های منابع محاسباتی هستند. الگوریتم‌های یادگیری تقویتی و بازی‌های دیفرانسیلی می‌توانند برای بهبود قابلیت‌های هدایت فضاپیماها، از جمله مقاومت در برابر تغییرات محیطی، کاهش تاخیرهای ناشی از ارتباطات زمینی، و افزایش کارایی محاسباتی، مورد استفاده قرار گیرند.


%در فرآیند طراحی مسیر، معمولاً یک مسیر بهینه و تاریخچه فرمان کنترلی طراحی می‌شود که با معیارهای ماموریت، مصرف سوخت و زمان پرواز، مطابقت داشته باشد. این روش قبل از پرواز انجام می‌شود و می‌تواند از استراتژی‌های متعددی برای هدایت بهینه با پیشران کم استفاده کند، از جمله تکنیک‌های بهینه‌سازی جهانی \cite{vavrina2017global} و برنامه‌نویسی غیرخطی \cite{ocampo2004finite}.
هدایت فضاپیماها معمولاً  پیش از پرواز انجام می‌شود. این روش‌ها می‌توانند از تکنیک‌های بهینه‌سازی فراگیر  \cite{vavrina2017global} یا برنامه‌نویسی غیرخطی برای تولید مسیرها و فرمان‌های کنترلی بهینه استفاده کنند. با این حال، این روش‌ها معمولا حجم محاسباتی زیادی دارند و برای استفاده درون‌سفینه نامناسب هستند \cite{ocampo2004finite}.
%با این حال، توانایی سریع محاسبه مجدد مسیر مرجع و تاریخچه فرمان کنترلی در داخل فضاپیما در هنگام پرواز بسیار مهم است. با نگاه به هدایت درون‌سفینه‌ای از دیدگاه یادگیری ماشین، یک کنترل‌کننده شبکه عصبی حلقه بسته، امکان محاسبه سریع و خودکار تاریخچه کنترل را برای یک فضاپیما فراهم می‌کند. به علاوه، یادگیری تقویتی\LTRfootnote{Reinforcement Learning (RL)}
% یادگیری پیش از پرواز واقعی است و می‌تواند از سخت‌افزارهای سریع و قدرتمند و ارزان قیمت زمینی جهت یادگیری استفاده کند.
% هدایت فضاپیماها معمولاً با استفاده از روش‌های طراحی مسیر پیش‌پرواز انجام می‌شود. با این حال، این روش‌ها محاسباتی فشرده هستند و برای استفاده درون‌سفینه نامناسب هستند.
 یادگیری ماشین می‌تواند برای بهبود قابلیت‌های هدایت فضاپیماها استفاده شود. کنترل‌کننده شبکه عصبی حلقه‌‌بسته می‌تواند برای محاسبه سریع و خودکار تاریخچه کنترل استفاده شود. یادگیری تقویتی نیز می‌تواند برای یادگیری رفتارهای هدایت بهینه استفاده شود.
 
روش‌های هدایت و بهینه‌سازی مسیر فضاپیماها به‌طور کلی به راه‌حل‌های اولیه مناسب نیاز دارند. در مسائل چند جسمی، طراحان مسیر اغلب حدس‌های اولیه کم‌هزینه‌ای برای انتقال‌ها با استفاده از نظریه سیستم‌های دینامیکی و منیفولدهای ثابت
\cite{2013AcAau, Haapala}
 ایجاد می‌کنند. 
%روش‌های مبتنی بر سیستم‌های دینامیکی در بسیاری از کاربردهای گذشته مفید بوده‌اند و در ترکیب با اصلاحات دیفرانسیل و یا تکنیک‌های بهینه‌سازی، برای بسیاری از کاربردها راه‌حل‌های بهینه تولید می‌کنند. با این حال، این رویکرد محاسباتی، فشرده است و اغلب به تعاملات انسان در حلقه نیاز دارد. به عنوان یک جایگزین، تکنیک‌های بهینه‌سازی جهانی ابتکاری مانند جستجوی حوضچه و الگوریتم‌های تکاملی، نیاز به راه‌حل‌های راه‌اندازی دقیق را کاهش می‌دهند 
%\cite{vavrina2017global}، اما پیچیدگی محاسباتی مربوطه آنها را برای استفاده درون‌سفینه ناممکن می‌سازد.

%روش‌های مبتنی بر سیستم‌های دینامیکی می‌توانند راه‌حل‌های بهینه تولید کنند، اما محاسباتی فشرده و نیازمند تعاملات انسان هستند. تکنیک‌های بهینه‌سازی جهانی می‌توانند نیاز به راه‌حل‌های راه‌اندازی دقیق \cite{vavrina2017global}
%را کاهش دهند، اما محاسباتی پیچیده هستند.



شبکه‌های عصبی ویژگی‌های جذابی برای فعال‌سازی هدایت در فضاپیما دارند. به‌عنوان مثال، شبکه‌های عصبی می‌توانند به‌طور مستقیم از تخمین‌های وضعیت به دستورهای پیشران کنترلی که با محدودیت‌های مأموریت سازگار است، برسند. عملکرد هدایت شبکه‌های عصبی در مطالعاتی مانند فرود بر سیارات \cite{gaudet2020six}، عملیات نزدیکی به سیارات \cite{gaudet2020terminal} و کنترل فضاپیما با پیشران ازدست‌رفته \cite{rubinsztejn2020neural} نشان داده شده است.
تازه‌ترین پیشرفت‌های تکنیک‌های یادگیری ماشین در مسائل خودکارسازی درونی به‌طور گسترده‌ای مورد مطالعه قرار گرفته‌اند؛ از پژوهش‌های اولیه تا توانایی‌های پیاده‌سازی.
به‌عنوان مثال، الگوریتم‌های یادگیری ماشین ابتدایی در فضاپیماهای مریخی نبرد برای کمک به شناسایی ویژگی‌های زمین‌شناسی تعبیه شده‌اند. الگوریتم \lr{AEGIS} توانایی انتخاب خودکار هدف توسط یک دوربین در داخل فضاپیماهای \lr{Spirit}، \lr{Opportunity} و \lr{Curiosity} را فعال دارد
\cite{estlin2012aegis}.
  در کامپیوتر پرواز اصلی، فرآیند دقت افزایی \lr{(Refinement Process)} نیاز به 94 تا 96 ثانیه دارد 
  \cite{francis2017aegis},
   که به طور قابل توجهی کمتر از زمان مورد نیاز برای ارسال تصاویر به زمین و انتظار برای انتخاب دستی توسط دانشمندان است.
   برنامه‌های آینده برای کاربردهای یادگیری ماشین درون‌سفینه شامل توانایی‌های رباتیکی درون‌سفینه برای فضاپیمای \lr{Perseverance}
\cite{higa2019vision, rothrock2016spoc}
     و شناسایی عیب برای 
   \lr{Europa Clipper}
\cite{wagstaff2019enabling} می‌شود. الگوریتم‌های یادگیری ماشین پتانسیلی برای سهم مهمی در مأموریت‌های اتوماسیون آینده دارند.


علاوه بر رباتیک سیاره‌ای، پژوهش‌های مختلفی به استفاده از تکنیک‌های مختلف یادگیری ماشین در مسائل نجومی پرداخته‌اند. در طراحی مسیر عملکرد رگرسیون معمولاً مؤثرتر هست. به عنوان مثال، از یک شبکه عصبی (\lr{NN}) در بهینه‌سازی مسیرهای رانشگر کم‌پیشران استفاده شده است 
\cite{dachwald2004evolutionary}.
 پژوهش‌های جدید شامل شناسایی انتقال‌های هتروکلینیک \cite{desmet2019identifying}، اصلاح مسیر رانشگر کم‌پیشران \cite{parrish2018lowthrust} و تجزیه و تحلیل مشکلات ازدست‌رفتن رانشگر \cite{rubinsztejn2020neural} می‌شود.
 
%تکنیک‌های یادگیری نظارتی می‌توانند نتایج مطلوبی تولید کنند، اما دارای محدودیت‌های قابل توجهی هستند. ابتدا، این رویکردها بر وجود دانش از پیش از فرآیند تصمیم‌گیری متکی هستند. کاربر با انتخاب نتایج مطلوب، فرض می‌کند که این دانش را دارد. این امر مستلزم دقیق بودن داده‌های تولید‌شده توسط کاربر برای نتایج مطلوب و همچنین وجود تکنیک‌های موجود برای حل مشکل کنونی و تولید داده است. در بخش‌هایی که چنین دانشی وجود ندارد، تکنیک‌های یادگیری نظارتی قابل استفاده نیستند.
تکنیک‌های یادگیری نظارتی می‌توانند نتایج مطلوبی تولید کنند؛ اما، دارای محدودیت‌های قابل توجهی هستند. یکی از این محدودیت‌ها این است که این رویکردها بر وجود دانش پیش از فرآیند تصمیم‌گیری متکی هستند. این امر مستلزم دقیق‌بودن داده‌های تولید‌شده توسط کاربر برای نتایج مطلوب و همچنین وجود تکنیک‌های موجود برای حل مشکل کنونی و تولید داده است.

%در سال‌های اخیر، \lr{RL} به اثبات مفید بودن خود در دستیابی به عملکرد بهترین حالت در دامنه‌هایی با ابهام محیطی قابل توجه رسیده است
%. هدایت فعال‌سازی‌شده توسط \lr{RL} به‌صورت گسترده‌ای بر اساس فاز پرواز دسته‌بندی می‌شوند. مسائل فرود \cite{furfaro2020adaptive, gaudet2020deep} 
%و عملیات در نزدیکی به اجسام کوچک
%\cite{gaudet2020terminal, gaudet2020six}
%از حوزه‌های مطالعاتی بهره‌برداری‌شده‌ای هستند. تحقیقات دیگر شامل مواجهه ،\cite{broida2019spacecraft} تداخل خارجی‌جوی
%\cite{gaudet2020reinforcement}
%، نگهداری ایستگاهی \cite{guzzetti2019reinforcement} و اجتناب از تشخیص \cite{reiter2020augmenting} هستند. مطالعاتی که فضاپیماهای رانشگر کم پیشرانرا در یک چارچوب دینامیکی چندبدنی با استفاده از \lr{RL} شامل طراحی انتقال با استفاده از
% \lr{Q-Learning}
% \cite{dasstuart2020rapid} و \lr{Proximal Policy Optimization }
%  \cite{miller2019lowthrust} و همچنین هدایت نزدیکی مدار
%  \cite{sullivan2020using}، شده‌اند.
  
  در سال‌های اخیر، قابلیت یادگیری تقویتی (\lr{RL}) در دستیابی به عملکرد بهینه در دامنه‌هایی با ابهام محیطی قابل توجه، به اثبات رسیده است  \cite{heess2017emergence, silver2017mastering}.
  هدایت انجام‌شده توسط \lr{RL} را می‌توان به‌صورت گسترده بر اساس فاز پرواز دسته‌بندی کرد.
  مسائل فرود \cite{furfaro2020adaptive, gaudet2020deep} 
   و عملیات در نزدیکی اجسام کوچک
   \cite{gaudet2020terminal, gaudet2020six}،
    از حوزه‌های پژوهشی هستند که از \lr{RL} استفاده می‌کنند.
  تحقیقات دیگر شامل مواجهه تداخل خارجی جوی \cite{gaudet2020reinforcement}،
   نگهداری ایستگاهی \cite{guzzetti2019reinforcement}  و هدایت به‌صورت جلوگیری از شناسایی \cite{reiter2020augmenting} است.
  مطالعاتی که فضاپیماهای رانشگر کم‌پیشران را در یک چارچوب دینامیکی چند بدنی با استفاده از RL انجام‌شده‌است، شامل طراحی انتقال با استفاده از \lr{Q-learning} 
   \cite{dasstuart2020rapid}، \lr{Proximal Policy Optimization}
     \cite{miller2019lowthrust}
      و  هدایت نزدیکی مدار  \cite{sullivan2020using} است.
      
      
      
      پژوهش‌های داخلی:
      
      
%  ناوبری و کنترل فضاپیماها و ماهواره‌ها، یکی از چالش‌های مهم در مهندسی فضایی است. این مسئله شامل تعیین موقعیت و جهت فضاپیما در فضا و هدایت و کنترل حرکت آن برای رسیدن به اهداف مورد نظر است.
%  در سال‌های اخ یر، پیشرفت‌های زیادی در زمینه ناوبری و کنترل فضاپیماها و ماهواره‌ها حاصل شده است. این پیشرفت‌ها عمدتاً ناشی از توسعه روش‌های جدیدی مانند یادگیری تقویتی، کنترل‌های تطبیقی و بینایی کامپیوتری است.
%  به عنوان مثال، در مرجع \cite{irani_1}، یک روش جدید برای انجام مانور اتصال دو فضاپیما با استفاده از الگوریتم یادگیری تقویتی  \lr{Q} ارائه شده است. این روش می‌تواند با دقت بسیار بالایی، مانور اتصال دو فضاپیما را انجام دهد.
%  در مرجع \cite{irani_2}، یک کنترل‌کننده تطبیقی بهینه بدون مدل برای بازپیکربندی آرایش پروازی ماهواره‌ها ارائه شده است. این کنترل‌کننده می‌تواند برای انواع مختلف آرایش پروازی ماهواره‌ها استفاده شود.
%  در مرجع \cite{irani_3}، یک روش جدید برای ناوبری و کنترل بازوی رباتیک ماهواره پرواز ـ آزاد ارائه شده است. این روش بر مبنای یک سیستم مبتنی بر بینایی کامپیوتری طراحی شده است.
%  پیشرفت‌های اخیر در زمینه ناوبری و کنترل فضاپیماها و ماهواره‌ها، به بهبود ایمنی و کارایی ماموریت‌های فضایی کمک کرده است.
  در زمینه هدایت و کنترل فضاپیماها و ماهواره‌ها، پیشرفت‌های زیادی در سال‌های اخیر حاصل شده است. این پیشرفت‌ها عمدتاً ناشی از توسعه روش‌های جدیدی مانند یادگیری تقویتی، کنترل‌های تطبیقی و بینایی کامپیوتر است. این پیشرفت‌ها به بهبود ایمنی و کارایی ماموریت‌های فضایی کمک کرده است.
  در مرجع \cite{irani_1}، یادگیری تقویتی برای مانور اتصال دو فضاپیما انجام شده است.
    در مرجع \cite{irani_2}، کنترل‌های تطبیقی برای بازپیکربندی آرایش پروازی ماهواره‌ها پیاده‌سازی شده است.
      در مرجع \cite{irani_3}، بینایی کامپیوتر برای ناوبری و کنترل بازوی رباتیک ماهواره اجرا شده است.
  

}{0}{{}}



\Description{اهداف پژوهش} {
\begin{itemize}
	\item  طراحی یک بازیکن جهت هدایت به کمک یادگیری تقویتی
	\item طراحی یک الگوریتم هدایت مقاوم به کمک بازی دیفرانسیلی
	\item بهبود عملکرد فضاپیما در محیط سه جسم
	\item عدم نیاز به ایستگاه زمینی جهت محاسبه صدور فرمان هدایت 
	\item پیاده‌سازی الگوریتم درون‌سفینه‌ای در محیط شبیه‌سازی‌شده

\end{itemize}}{0}{{}}

\vspace{-1.5cm}
\renewcommand{\arraystretch}{1}
\begin{table}[H]
	\centering
	\caption{جدول برنامه‌ریزی پروژه ناوبری اینرسی مشارکتی مبتنی بر هوش مصنوعی}
	\vspace{.2cm}
	\begin{tabular}{|P{3cm}|P{10cm}|}
		\hline
		\textbf{فصل} & \textbf{فعالیت‌ها} \\
		\hline
		\multirow{4}{*}{\textbf{تابستان سال اول}} & انجام مطالعه و مروری جامع در زمینه بازی دیفرانسیلی و یادگیری تقویتی.\\
		\cline{2-2}
		&
        بررسی و تعریف دقیق مسأله یادگیری تقویتی و الگوریتم هدایت به وسیله یادگیری تقویتی در مسئله سه‌جسم.
          \\
		\cline{2-2}
		& طراحی ساختار محیط سه‌جسم و بازیکن یادگیری تقویتی و ارتباط بین محیط و بازیکن. \\
		\hline
		\multirow{3}{*}{\textbf{نیمسال اول سال دوم}} & طراحی مسیر انتقال بهینه بین دو مدار \\
		\cline{2-2}
		& شبیه‌سازی دقیق محیط مسأله سه‌جسم جهت آموزش بازیکن یادگیری تقویتی \\
		\cline{2-2}
		&  طراحی و توسعه سیستم یادگیری تقویتی جهت محاسبه دستور \\
		\hline
		\multirow{2}{*}{\textbf{نیمسال دوم سال دوم}} & پیاده‌سازی هدایت فضاپیما با استفاده از الگوریتم‌های کلاسیک \\
		\cline{2-2}
		& آموزش الگوریتم‌های یادگیری تقویتی در محیط‌های شبیه‌سازی شده\\
		\hline
		\multirow{2}{*}{\textbf{تابستان سال دوم}}  
		&  اصلاح و بهبود عملکرد یادگیری تقویتی در هدایت فضاپیما \\ 
		\cline{2-2}
		\cline{2-2}
		& نتیجه‌گیری از تحقیقات و ارائه گزارش نهایی شامل معرفی مسئله، روش‌های استفاده‌شده، نتایج حاصل و پیشنهادات برای تحقیقات آتی \\
		\hline
	\end{tabular}
\end{table}





\Description{نتایج مورد انتظار} {
	\begin{itemize}
				\item  افزایش ایمنی ماموریت به‌دلیل استفاده از بازی دیفرانسیلی
		\item  کاهش مصرف سوخت و هزینه به‌دلیل بهینه‌بودن الگوریتم هدایت
		\item  هدایت فضاپیما بدون نیاز به ایستگاه زمینی
		\item کاهش هزینه ساخت به‌دلیل عدم نیاز به سیستم‌های ارتباطی قوی
		
\end{itemize}}{0}{{}}
\vspace{-1cm}
\Description{روش صحت‌سنجی نتایج} {
	\begin{itemize}
		\item  مقایسه با سایر روش‌های معتبر
		\item مدل‌سازی و شبیه‌سازی سیستم محیط و بازیکن
		
\end{itemize}}{0}{{}}
\vspace{-1cm}
\Description{گلوگاه‌های پیش‌بینی‌‌شده} {
	\begin{itemize}
		\item  آموزش شبکه: به‌علت پیچیده‌بودن شبکه و محیطی که طراحی شده است، محاسبات سنگین می‌شود.
		\item پایداری و کنترل: به‌علت غیرخطی‌بودن دینامیک سامانه و ناپایداری ذاتی سیستم، کنترل و پایداری آن مشکل است.
		\vspace{-0.9cm}
\end{itemize}}{0}{{}}


\Description{‌‌نوآوری‌ها} {
%	\begin{itemize}
%		\item  طراحی الگوریتم هدایت بهینه درون‌سفینه‌ای
%		\item  طراحی الگوریتم هدایت مقاوم درون‌سفینه‌ای
%		\item طراحی الگوریتم هدایت یادگیری تقویتی بر مبنای بازی دیفرانسیلی
%		
%\end{itemize}
پیش از این، از ترکیب یادگیری تقویتی و بازی دیفرانسیلی در یک مسئله چندجسمی پویا استفاده نشده است. به‌علاوه، منطق هدایت ساخته‌شده از ترکیب دو روش، یک هدایت مقاوم و در عین حال بهینه  است، که تاکنون در این مسئله انجام نشده است.
}{0}{{}}



%\bibliography{refs.bib}
%\input{refs.bib}
%%\bibliography{refs.bib}
%\bibliographystyle{unsrt}


%\printbibliography

%\bibliography{refs}


\noindent
\newpage

% Add our references in references,bib
%\References{}
%
% -------------------------------------------------------
%  Bibliography
% -------------------------------------------------------


\begin{latin}
%	\renewcommand{\bibname}{\rl{مراجع}}
\baselineskip=.8\baselineskip
\bibliographystyle{./styles/packages/unsrtabbrv}

% Uncomment next line to include uncited references
% \nocite{*}

\bibliography{bibs/full,bibs/refs}
 
\end{latin}
\newpage


%\begin{thebibliography}{1}
%	
% -------------------------------------------------------
%  Bibliography
% -------------------------------------------------------


\begin{latin}
%	\renewcommand{\bibname}{\rl{مراجع}}
\baselineskip=.8\baselineskip
\bibliographystyle{./styles/packages/unsrtabbrv}

% Uncomment next line to include uncited references
% \nocite{*}

\bibliography{bibs/full,bibs/refs}
 
\end{latin}
\newpage

%\end{thebibliography}

\begin{latin}
	\begin{thebibliography}{1}
		
		
		\bibitem{vavrina2017global}
		Vavrina, Matthew A, Jacob A Englander, Sean M Phillips, and Kyle M Hughes.
		"Global, multi-objective trajectory optimization with parametric spreading."
		In AAS AIAA Astrodynamics Specialist Conference 2017, Tech. No. GSFC-E-DAA-TN45282, 2017.	
		
		
		\bibitem{ocampo2004finite}
		Ocampo, Cesar.
		"Finite Burn Maneuver Modeling for a Generalized Spacecraft Trajectory Design and Optimization System."
		Annals of the New York Academy of Sciences 1017 (2004): 210-233.
		doi: 10.1196/annals.1311.013.
		
		
%		\bibitem{Haapala}
%		Haapala, A. F., \& Howell, K. C. (2016).
%		"A framework for constructing transfers linking periodic libration point orbits in the spatial circular restricted three-body problem."
%		International Journal of Bifurcation and Chaos, 26(05), 1630013.
		
		
		\bibitem{2013AcAau}
		Marchand, B.~G. and Scarritt, S.~K. and Pavlak, T.~A. and Howell, K.~C.
		``A dynamical approach to precision entry in multi-body regimes: Dispersion manifolds.''
		Acta Astronautica 89 (2013): 107-120.
		doi: 10.1016/j.actaastro.2013.02.015.
		
		\bibitem{estlin2012aegis}
		T.A. Estlin, B.J. Bornstein, D.M. Gaines, R.C. Anderson, D.R. Thompson, M. Burl, R. Castaño, M. Judd.
		"Aegis automated science targeting for the MER Opportunity rover."
		ACM Trans. Intell. Syst. Technol. (TIST), 3, 1-19, 2012.
		
		\bibitem{francis2017aegis}
		R. Francis, T. Estlin, G. Doran, S. Johnstone, D. Gaines, V. Verma, M. Burl, J. Frydenvang, S. Montano, R. Wiens, S. Schaffer, O. Gasnault, L. Deflores, D. Blaney, B. Bornstein.
		"Aegis autonomous targeting for ChemCam on Mars Science Laboratory: Deployment and results of initial science team use."
		Science Robotics, 2, 2017.
		
		\bibitem{wagstaff2019enabling}
		K.L. Wagstaff, G. Doran, A. Davies, S. Anwar, S. Chakraborty, M. Cameron, I. Daubar, C. Phillips.
		"Enabling onboard detection of events of scientific interest for the Europa Clipper spacecraft,"
		in: 25th ACM SIGKDD International Conference on Knowledge Discovery \& Data Mining,
		Association for Computing Machinery, Anchorage, Alaska, 2019, pp. 2191--2201,
		doi: 10.1145/3292500.3330656.
		
		\bibitem{higa2019vision}
		S. Higa, Y. Iwashita, K. Otsu, M. Ono, O. Lamarre, A. Didier, M. Hoffmann.
		"Vision-based estimation of driving energy for planetary rovers using deep learning and terramechanics,"
		IEEE Robot. Autom. Lett. 4 (2019) 3876–3883.
		
		\bibitem{rothrock2016spoc}
		B. Rothrock, J. Papon, R. Kennedy, M. Ono, M. Heverly, C. Cunningham.
		"Spoc: Deep learning-based terrain classification for Mars rover missions,"
		in: AIAA Space and Astronautics Forum and Exposition, SPACE 2016,
		American Institute of Aeronautics and Astronautics Inc, AIAA, 2016, pp. 1–12.
		
		\bibitem{dachwald2004evolutionary}
		B. Dachwald.
		"Evolutionary neurocontrol: A smart method for global optimization of low-thrust trajectories,"
		in: AIAA/AAS Astrodynamics Specialist Conference and Exhibit,
		Providence, Rhode Island, 2004, pp. 1–16.
		
		\bibitem{desmet2019identifying}
		S. De Smet, D.J. Scheeres.
		"Identifying heteroclinic connections using artificial neural networks,"
		Acta Astronaut. 161 (2019) 192–199.
		
		\bibitem{parrish2018lowthrust}
		N.L.O. Parrish.
		"Low Thrust Trajectory Optimization in Cislunar and Translunar Space (Ph.D. thesis),"
		University of Colorado Boulder, 2018.
		
		\bibitem{rubinsztejn2020neural}
		A. Rubinsztejn, R. Sood, F.E. Laipert.
		"Neural network optimal control in astrodynamics: Application to the missed thrust problem,"
		Acta Astronaut. 176 (2020) 192–203.
		
		\bibitem{heess2017emergence}
		N. Heess, D. TB, S. Sriram, J. Lemmon, J. Merel, G. Wayne, Y. Tassa, T. Erez, Z. Wang, S.M.A. Eslami, M.A. Riedmiller, D. Silver.
		"Emergence of locomotion behaviours in rich environments," 2017.
		CoRR abs/1707.02286.
		arXiv:1707.02286.
		
		\bibitem{silver2017mastering}
		D. Silver, J. Schrittwieser, K. Simonyan, I. Antonoglou, A. Huang, A. Guez, T. Hubert, L. Baker, M. Lai, A. Bolton, Y. Chen, T. Lillicrap, F. Hui, L. Sifre, G. van den Driessche, T. Graepel, D. Hassabis.
		"Mastering the game of Go without human knowledge," Nature 550 (2017) Article.
		
		\bibitem{furfaro2020adaptive}
		R. Furfaro, A. Scorsoglio, R. Linares, M. Massari,
		"Adaptive generalized ZEM-ZEV feedback guidance for planetary landing via a deep reinforcement learning approach,"
		Acta Astronaut. 171 (2020) 156–171.
		
		\bibitem{gaudet2020deep}
		B. Gaudet, R. Linares, R. Furfaro,
		"Deep reinforcement learning for six degrees of freedom planetary landing,"
		Adv. Space Res. 65 (2020) 1723–1741.
		
        \bibitem{broida2019spacecraft}
		J. Broida, R. Linares,
		"Spacecraft rendezvous guidance in cluttered environments via reinforcement learning,"
		in: \textit{29th AAS/AIAA Space Flight Mechanics Meeting}, American Astronautical Society, Ka’anapali, Hawaii, 2019, pp. 1–15.
		
		\bibitem{gaudet2020reinforcement}
		B. Gaudet, R. Furfaro, R. Linares,
		"Reinforcement learning for angle-only intercept guidance of maneuvering targets,"
		\textit{Aerosp. Sci. Technol.} 99 (2020).
		
		\bibitem{guzzetti2019reinforcement}
		D. Guzzetti,
		"Reinforcement learning and topology of orbit manifolds for station-keeping of unstable symmetric periodic orbits,"
		in: \textit{AAS/AIAA Astrodynamics Specialist Conference}, American Astronautical Society, Portland, Maine, 2019, pp. 1–20.
		
		\bibitem{reiter2020augmenting}
		J.A. Reiter, D.B. Spencer,
		"Augmenting spacecraft maneuver strategy optimization for detection avoidance with competitive coevolution,"
		in: \textit{20th AIAA Scitech Forum}, AIAA, Orlando, Florida, 2020, pp. 1–11.
		
		\bibitem{dasstuart2020rapid}
		A. Das-Stuart, K.C. Howell, D.C. Folta,
		"Rapid trajectory design in complex environments enabled by reinforcement learning and graph search strategies,"
		\textit{Acta Astronaut.} 171 (2020) 172–195.
		
		\bibitem{miller2019lowthrust}
		D. Miller, R. Linares,
		"Low-thrust optimal control via reinforcement learning,"
		in: \textit{29th AAS/AIAA Space Flight Mechanics Meeting}, American Astronautical Society, Kaanapali, Hawaii, 2019, pp. 1–18.
		
		\bibitem{sullivan2020using}
		C.J. Sullivan, N. Bosanac,
		"Using reinforcement learning to design a low-thrust approach into a periodic orbit in a multi-body system,"
		in: \textit{20th AIAA Scitech Forum}, AIAA, Orlando, Florida, 2020, pp. 1–19.
		
		
		\bibitem{gaudet2020terminal}
		B. Gaudet, R. Linares, R. Furfaro,
		"Terminal adaptive guidance via reinforcement meta-learning: Applications to autonomous asteroid close-proximity operations,"
		\textit{Acta Astronaut.} 171 (2020) 1–13.
		
		
		\bibitem{gaudet2020six}
		B. Gaudet, R. Linares, R. Furfaro,
		"Six degree-of-freedom hovering over an asteroid with unknown environmental dynamics via reinforcement learning,"
		in: \textit{20th AIAA Scitech Forum}, AIAA, Orlando, Florida, 2020, pp. 1–15.
		
		\bibitem{irani_1}
		صیامی عراقی, م.,  شفیعی نژاد, ا. (1399). طراحی مسیر اتصال میان دو ماهواره با الگوریتم یادگیری تقویتی Q. پنجمین کنفرانس ملی مهندسی مکانیک و هوافضا. https://civilica.com/doc/1039632
		
				\bibitem{irani_2}
کنکاشور, محمدرسول, بلندی, حسین,  مزینی, ناصر. (1401). طراحی کنترل‌کننده تطبیقی بهینه بدون مدل برای باز پیکربندی آرایش پروازی ماهواره‌ها با یادگیری تقویتی. دانش و فناوری هوافضا, 11(1), 41-60.
		
				\bibitem{irani_3}
		موسوی, فرزاد (1401). ناوبری و کنترل بازوی رباتیک ماهواره پرواز ـ آزاد حین جا به‌ جایی و قراردادن محموله در ایستگاه‌ فضایی. http://library.sharif.ir/parvan/resource/496295
		
	\end{thebibliography}
\end{latin}

% Change \newcommand{\Courses} in the previous lines
%\Courses{}

%\noindent

%\Signatures{}

\end{document}
